\documentclass[11pt]{article}

\usepackage[margin=1in]{geometry}
\usepackage{amsmath,amssymb,amsthm,mathtools}
\usepackage{graphicx}
\usepackage{xcolor}
\usepackage{microtype}
\usepackage[colorlinks=true,linkcolor=blue!55!black,citecolor=blue!55!black,urlcolor=blue!55!black]{hyperref}

\newtheorem{theorem}{Theorem}
\newtheorem{lemma}[theorem]{Lemma}
\newtheorem{remark}[theorem]{Remark}
\newcommand{\Rad}{\operatorname{Rad}}
\newcommand{\Alg}{\operatorname{Alg}}

\title{A power-compact weighted-shift subcase\\
of the positive-commutator radical problem}
\author{Automated mathematical research packet\\\small Run fa\_banach\_001, lane 0}
\date{21 July 2026}

\begin{document}
\maketitle

\begin{abstract}
Kandi\'c and \v{S}ivic asked whether their theorem for two positive
power-compact operators remains true when only one factor is power compact:
must some power of a positive commutator belong to the radical of the Banach
algebra generated by the two factors?  We prove an affirmative subcase.  If
the power-compact factor is a unilateral forward weighted shift with strictly
positive weights on $c_0$ or $\ell^p$, then the $3m$-th power of the
commutator is radical whenever the $m$-th power of the shift is compact.  The
key new observation is that order semicommutation forces the other positive
operator to preserve the shift's maximal tail nest.
\end{abstract}

\section{The source question}

Let $E$ be a Banach lattice, let $A,B\in\mathcal L(E)$ be positive, and put
$C=AB-BA\geq0$.  The source paper proves that if both $A$ and $B$ are power
compact, then some power of $C$ belongs to the radical of the Banach algebra
generated by $A$ and $B$ \cite[Theorem~4.7]{KS}.  It ends by asking whether
only one of $A,B$ need be power compact.

\begin{center}
\includegraphics[width=.94\linewidth]{figures/open_problem_crop.png}
\end{center}

The earlier theorem of Drnov\v{s}ek already shows that if only one factor is
power compact, then $C$ is power compact, ideal-triangularizable, and
quasinilpotent \cite[Corollary~3.2]{D}.  What is missing in general is control
of products $DC^n$ for arbitrary $D$ in the algebra generated by $A$ and
$B$.

\section{Statement of the partial result}

Let $E$ be either $c_0(\mathbb N_0)$ or $\ell^p(\mathbb N_0)$,
$1\leq p<\infty$, over $\mathbb R$ or $\mathbb C$.  Write $(e_j)_{j\geq0}$
for its standard unit vectors.

\begin{theorem}\label{thm:main}
Let $A$ be the positive unilateral forward weighted shift
\[
       Ae_j=\alpha_j e_{j+1}\qquad(j\geq0),
       \qquad \alpha_j>0.
\]
Assume that $A^m$ is compact for some $m\geq1$.  If $B\geq0$ and
$AB\geq BA$, set $C=AB-BA$ and let $\mathcal A=\Alg(A,B)$ be the norm-closed
Banach algebra generated by $A$ and $B$.  Then
\[
                 C^{3m}\in\Rad(\mathcal A).
\]
\end{theorem}

The exponent is explicit but is not claimed to be optimal.

\section{The forced tail nest}

For $n\geq0$ define the closed tail ideal
\[
        M_n=\overline{\operatorname{span}}\{e_j:j\geq n\}.
\]
Together with $\{0\}$, these spaces form a complete maximal chain: successive
quotients $M_n/M_{n+1}$ are one dimensional and $\bigcap_nM_n=\{0\}$.

\begin{lemma}[boundary propagation]\label{lem:tail}
Under the hypotheses of Theorem~\ref{thm:main}, every $M_n$ is invariant
under $B$.  Consequently $A$, $B$, and every member of $\mathcal A$ are
triangular with respect to the tail chain.  The diagonal coefficients of
$C$ on this chain are all zero.
\end{lemma}

\begin{proof}
Write
\[
       b_{ij}=e_i^*(Be_j)\geq0.
\]
The $(i,j)$ matrix coefficient of $AB\geq BA$ says
\begin{align}
 0&\geq \alpha_j b_{0,j+1} &&(i=0),\label{eq:boundary}\\
 \alpha_{i-1}b_{i-1,j}&\geq \alpha_jb_{i,j+1} &&(i\geq1).
 \label{eq:recursion}
\end{align}
All weights are strictly positive.  Equation \eqref{eq:boundary} therefore
gives $b_{0k}=0$ for every $k>0$.  Suppose inductively that
$b_{i-1,k}=0$ whenever $k>i-1$.  If $k>i$, apply
\eqref{eq:recursion} with $j=k-1$; its left side is zero, so $b_{ik}=0$.
Thus
\[
                         b_{ik}=0\qquad(k>i).
\]
Equivalently, $B(M_n)\subseteq M_n$ for every $n$.  The same is clear for
$A$, and tail invariance is preserved by algebraic operations and norm
limits.

Finally, for $n\geq1$,
\[
 (AB)_{nn}=\alpha_{n-1}b_{n-1,n}=0,
 \qquad
 (BA)_{nn}=\alpha_nb_{n,n+1}=0;
\]
the $n=0$ case follows from \eqref{eq:boundary}.  Hence every diagonal
coefficient of $C$ is zero.
\end{proof}

\section{A compact power and the radical criterion}

We use the standard Aliprantis--Burkinshaw domination theorem: if
$0\leq S\leq T$ on a Banach lattice and $T$ is compact, then $S^3$ is
compact \cite[Theorem~5.14]{AB}.  This is the same domination result used in
the proof of \cite[Corollary~3.2]{D}.

\begin{lemma}\label{lem:compact}
Under the hypotheses of Theorem~\ref{thm:main}, $C^{3m}$ is compact.
\end{lemma}

\begin{proof}
Because $0\leq C\leq AB$, positivity and induction give
$0\leq C^m\leq(AB)^m$.  The inequality $BA\leq AB$ similarly permits each
$B$ to be moved to the right, giving
\[
                         (AB)^m\leq A^mB^m.
\]
The positive operator $A^mB^m$ is compact.  Thus
$0\leq C^m\leq A^mB^m$, and the domination theorem yields
$(C^m)^3=C^{3m}$ compact.
\end{proof}

\begin{proof}[Proof of Theorem~\ref{thm:main}]
Pass to the standard complexification in the real case.  Let
$D\in\mathcal A^\#$, the unitization of $\mathcal A$.  Lemma~\ref{lem:tail}
shows that $D$ and $C$ are triangular on the maximal tail chain, and the
diagonal of $C$ is zero.  Therefore every diagonal coefficient of
$DC^{3m}$ is zero.  By Lemma~\ref{lem:compact}, this operator is compact.
Ringrose's theorem for compact operators and complete triangularizing chains
now gives
\[
                         \sigma(DC^{3m})=\{0\}.
\]
Thus $r(DC^{3m})=0$ for every $D\in\mathcal A^\#$.  The spectral-radius
characterization of the Jacobson radical proves
$C^{3m}\in\Rad(\mathcal A)$.
\end{proof}

\section{Scope and obstruction to the full result}

The proof separates the source problem into two issues.  Power compactness of
the commutator is not the main obstruction: from one compact power of $A$,
order domination already supplies the compact power $C^{3m}$.  The difficult
point is to make every product $DC^{3m}$ quasinilpotent.  For a forward
weighted shift, the boundary row in \eqref{eq:boundary} propagates through
the matrix and creates a common maximal nest.  A general positive
power-compact operator has no analogous coordinate boundary, so the argument
does not settle the full question.

If the power-compact factor is actually nilpotent, say $A^m=0$, the same
inequalities give $0\leq C^m\leq A^mB^m=0$, and the source conclusion is
trivial.  The theorem above is useful precisely when $A$ is power compact but
not nilpotent.

\begin{thebibliography}{9}

\bibitem{AB}
C.~D. Aliprantis and O.~Burkinshaw,
\emph{Positive Operators}, Springer, 2006.

\bibitem{D}
R.~Drnov\v{s}ek,
Once more on positive commutators,
\emph{Studia Math.} \textbf{211} (2012), 241--245;
\href{https://arxiv.org/abs/1211.0812}{arXiv:1211.0812}.

\bibitem{KS}
M.~Kandi\'c and K.~\v{S}ivic,
On the positive commutator in the radical,
\emph{Positivity} \textbf{21} (2017), 99--111;
\href{https://arxiv.org/abs/1405.4988}{arXiv:1405.4988}.

\end{thebibliography}

\end{document}
